% =====================================================================
% Chapter 2: Literature Review
% =====================================================================
\chapter{Literature Review}
\label{ch:literature}

\section{Introduction}

Jersey number recognition sits at the intersection of three broad
research streams: object detection and pose estimation, scene text
recognition and OCR, and sports video analytics. Each of these areas
has matured substantially over the past decade, and the specific design
choices made in this thesis --- YOLOv8x with pose-driven torso
localization for images, a football-specific YOLOv5 detector with
ByteTrack for video, a multi-variant EasyOCR pipeline, and a
transformer-based PARSeq recognizer with an EasyOCR fallback --- are
each grounded in this prior work. This chapter surveys the most
relevant developments in each area, discusses their reported strengths
and limitations, and explains how they informed the design decisions
taken in this thesis.

\section{Object Detection and Pose Estimation}

Object detection research has been shaped by the shift from two-stage,
region-proposal-based detectors toward single-stage detectors that
frame localization and classification as a single dense regression
problem. The YOLO (You Only Look Once) family pioneered this
single-stage paradigm, trading a small amount of accuracy for a very
large gain in inference speed, which is essential for processing the
thousands of player crops and video frames evaluated in this thesis
\cite{redmon2016yolo}. The YOLOv8 architecture used in the image-based
pipeline of this thesis extends its predecessors with an anchor-free
detection head, decoupled classification and regression branches, and
refined training recipes, and it is available pretrained on large-scale
datasets such as COCO and ImageNet, allowing this thesis to use it
directly as a generic person detector without task-specific
fine-tuning \cite{ultralytics2023yolov8,deng2009imagenet}. For the
video-based pipeline, a YOLOv5 model fine-tuned specifically on
football broadcast footage is used instead of a generic person
detector; this specialist model detects four semantic classes ---
player, goalkeeper, referee, and ball --- and consistently outperforms
generic person detectors in football-specific scenes because it has
been trained on data that already reflects the distinctive appearance
of players, kits, and camera angles found in broadcast football.

Human pose estimation provides a complementary signal to plain bounding
boxes: rather than treating a detected player as an undifferentiated
rectangle, pose keypoints allow specific anatomical landmarks to be
located within it. Part-affinity-field approaches such as OpenPose
demonstrated that multiple people could be jointly and efficiently
localized with a single forward pass, associating individual keypoints
into complete skeletons even in crowded scenes \cite{cao2017openpose}.
In the image pipeline of this thesis, YOLOv8x-Pose provides 17
COCO-style human keypoints per detected person. Shoulder and hip
keypoints are used to derive a tight torso bounding box, substantially
narrowing the region that must be searched for a jersey number compared
with using the full player box, and reducing the amount of background,
limbs, and shorts that would otherwise dilute the OCR region of
interest.

\section{Multi-Object Tracking}

Reliable player tracking across video frames is essential for the
video-based pipeline, since jersey-number evidence for the same
individual must be accumulated over multiple frames rather than
inferred from a single, possibly blurred, observation. Classical
tracking-by-detection approaches such as SORT associate detections
across frames purely through Kalman-filter motion prediction and
Hungarian-algorithm bounding-box matching, offering a simple and fast
baseline but suffering from frequent identity switches whenever
detections are momentarily missed \cite{bewley2016sort}. DeepSORT
extended this approach with a learned appearance descriptor, reducing
identity switches at the cost of an additional embedding network that
must be run for every detection \cite{wojke2017deepsort}. ByteTrack
takes a different approach: rather than discarding low-confidence
detections outright, it performs a secondary association step using
these low-confidence boxes to recover otherwise lost tracks, achieving
state-of-the-art tracking accuracy while remaining detection-only and
therefore computationally efficient \cite{zhang2022bytetrack}.
ByteTrack does not require a separate re-identification network, making
it well-suited to integration with detection-only models such as the
YOLOv5 football detector used in this thesis, where every extra
per-frame inference call directly adds to the already substantial OCR
processing cost analyzed in Chapter~\ref{ch:results}.

\section{Scene Text Recognition and OCR for Jersey Numbers}

General-purpose OCR engines such as Tesseract were designed for
high-contrast, well-segmented document text and consistently fail on
scene text --- text appearing in natural images under unconstrained
illumination, orientation, and background clutter. The shift from
document OCR to scene text recognition was driven substantially by
end-to-end trainable sequence models: CRNN combines a convolutional
feature extractor with a recurrent sequence model and a CTC decoding
layer, removing the need for explicit character segmentation
\cite{shi2017crnn}. Robust scene-text \emph{detection}, as opposed to
recognition, was advanced by CRAFT, which localizes individual
character regions and their affinity to neighboring characters rather
than predicting a single bounding box per word, making it considerably
more robust to the curved, rotated, and partially occluded text
encountered on jersey fabric \cite{baek2019craft}. EasyOCR, used
extensively in the image-based pipeline of this thesis, combines a
CRAFT-style text detector with a CRNN-style recognizer and provides a
practical, freely available starting point for digit recognition in
natural images; however, as documented later in
Chapter~\ref{ch:methodology}, its out-of-the-box performance on small,
low-resolution jersey-digit regions remains insufficient without the
multi-variant enhancement, zoom, and shape-disambiguation strategy
developed in this thesis.

More recently, transformer-based sequence models have overtaken
CRNN-style recognizers on standard scene-text benchmarks. PARSeq
formulates scene text recognition as a permuted autoregressive sequence
modeling problem, allowing the same architecture to be trained and
used in both autoregressive and non-autoregressive decoding modes and
achieving strong accuracy on curved and occluded text while remaining
attention-based rather than recurrent \cite{bautista2022parseq}, an
approach that builds directly on the general-purpose Transformer
architecture \cite{vaswani2017attention}. The video-based pipeline in
this thesis adopts PARSeq as its primary OCR backend for jersey digits,
attempting to load the pretrained PARSeq model at start-up and falling
back automatically to EasyOCR whenever PARSeq cannot be loaded in the
runtime environment, which provides a practical balance between the
stronger sequence modeling of a transformer-based recognizer and the
robustness of a widely deployed fallback engine.

Several prior works have specifically targeted jersey number recognition
as a component of player re-identification in sports video
\cite{gerke2015soccer,li2018jersey}. Gerke et al.\ combine a
hand-crafted HOG feature representation with an SVM classifier and a
temporal voting stage over multiple frames of the same player, an
architecture that pre-dates the deep-learning era but already
identifies the two design principles that recur throughout the later
literature and throughout this thesis: (a) accurate localization of the
digit region is the dominant factor in recognition accuracy, and (b)
temporal aggregation of single-frame results across multiple
observations of the same player substantially improves the final
prediction \cite{gerke2015soccer}. Li et al.\ replace the hand-crafted
pipeline with a CNN-LSTM architecture and add an explicit legibility
filter so that OCR is only attempted on crops likely to contain a
readable digit, reporting a further accuracy gain from this filtering
step \cite{li2018jersey}. Both observations are directly reflected in
the design of the pipelines presented in this thesis: pose-driven torso
localization narrows the OCR search region in the image pipeline, a
temporal jersey voting mechanism aggregates evidence in the video
pipeline, and a dedicated ResNet-18 legibility classifier filters
unreadable crops in the image pipeline.

\section{The SoccerNet Dataset}

The SoccerNet initiative provides one of the largest collections of
annotated broadcast football video, with tasks spanning action
spotting, player tracking, camera calibration, and jersey number
recognition \cite{deliege2021soccernet}. The SN-Jersey-2023 sub-task
provides cropped player images organized by player identity, with
ground-truth jersey numbers in JSON format; this thesis uses a curated
subset of this sub-task for the image-based pipeline. The SoccerNet
Tracking sub-task, together with videos collected from the Hugging Face
SoccerNet SN-PCBAS release and from Rutracker.org, provides
full-resolution broadcast clips with per-frame player bounding-box
annotations in MOT format, which serve as the source of both the ten
video clips and the tracking ground-truth files used in the video
pipeline of this thesis. Because SoccerNet does not provide manually
verified jersey-number ground truth for these particular tracking
clips, this thesis additionally produced manually annotated
jersey-number CSV files using CVAT, as described in
Chapter~\ref{ch:methodology}, so that video-pipeline accuracy could be
reported against genuine ground truth rather than against the
pipeline's own internal predictions.

\section{Legibility Classification}

Not every cropped player image contains a recognizable jersey number.
Players may face the camera directly, showing only a front sponsor
logo rather than a jersey number, be heavily occluded by another
player or the ball, or appear at a resolution too low to resolve the
digits at all. Several prior jersey recognition works include a
legibility pre-filter so that OCR is only attempted on crops that are
likely to yield a valid result \cite{li2018jersey}. Training such a
filter typically requires either costly manual annotation of legibility
labels or, as in this thesis, an automatic weak-labeling strategy in
which a crop is labeled legible whenever the multi-variant OCR pipeline
already returns a sufficiently confident digit reading. This thesis
trains a ResNet-18 binary classifier for this purpose, using the
ImageNet-pretrained backbone popularized by residual networks
\cite{he2016resnet,deng2009imagenet}, thereby avoiding the need for
costly manual annotation of legibility while still obtaining a learned,
image-based signal that is complementary to the OCR engine's own
confidence score.

\section{Summary}

The works reviewed in this chapter collectively motivate the principal
design choices of this thesis: YOLOv8x for image-based player
detection due to its strong accuracy and available pose variant; a
specialist YOLOv5 football model for video-based detection, providing
semantic class labels for player, goalkeeper, and referee; ByteTrack
for multi-frame identity association in video, chosen over SORT and
DeepSORT for its favorable accuracy-to-cost trade-off; a multi-variant
EasyOCR pipeline with domain-specific enhancement, shape-based digit
disambiguation, and dynamic zoom to compensate for the well-documented
limitations of general-purpose OCR on jersey digits in the image
pipeline; a PARSeq transformer-based recognizer with an EasyOCR
fallback for the video pipeline, reflecting the shift toward
attention-based scene text recognition; and a ResNet-18 legibility
classifier together with a temporal jersey voting mechanism to provide
complementary signals for improving recognition reliability.
Chapter~\ref{ch:methodology} describes how these components were
implemented and combined.
